\documentclass[10pt]{article}
\usepackage[T1]{fontenc}
\usepackage[utf8]{inputenc}
\usepackage[french]{babel}
\usepackage{amsmath,amssymb,mathrsfs,mathpazo}
\usepackage{geometry}
\geometry{a4paper,margin=2cm}
\usepackage{graphicx,array,tabularx,tabularray,colortbl,multirow,multicol,booktabs,enumitem,pifont,ulem,fancybox,diagbox,cancel}
\usepackage{pstricks,pstricks-add,pst-plot,pst-tree,pst-node,pst-text,pst-func,pst-math,pst-bspline,pst-3dplot,pst-eucl,pst-solides3d}
\let\pscalculate\undefined
\usepackage{pst-bezier}
\usepackage{tikz,pgfplots,tkz-tab}
\usetikzlibrary{arrows,arrows.meta,patterns,calc,positioning,trees,decorations.markings,shapes.geometric,shapes.misc,shapes.arrows,intersections,decorations.pathmorphing,decorations.pathreplacing,backgrounds,shadows,babel}
\pgfplotsset{compat=1.18}
\usepgfplotslibrary{fillbetween}
\usepackage[np]{numprint}
\usepackage{listings,makecell,pas-tableur,twoopt}
\let\barre\undefined
\usepackage[pstricks]{tablvar}
\definecolor{darkgreen}{RGB}{0,100,0}
\providecommand{\red}{\color{red}}
\providecommand{\blue}{\color{blue}}
\providecommand{\green}{\color{green}}
\providecommand{\black}{\color{black}}
\providecommand{\white}{\color{white}}
\providecommand{\gray}{\color{gray}}
\providecommand{\cyan}{\color{cyan}}
\providecommand{\magenta}{\color{magenta}}
\providecommand{\yellow}{\color{yellow}}
\providecommand{\violet}{\color{violet}}
\providecommand{\orange}{\color{orange}}
\providecommand{\euro}{\texteuro}
\providecommand{\EUR}{\texteuro}
\providecommand{\ieme}{\textsuperscript{e}}
\providecommand{\degres}{\ensuremath{^\circ}}
\providecommand{\degree}{\ensuremath{^\circ}}
\providecommand{\vv}[1]{\overrightarrow{#1}}
\providecommand{\Coord}[1]{#1}
\providecommand{\decosix}{}
\providecommand{\starredbullet}{\ensuremath{\bullet}}
\pagestyle{empty}
\setlength{\parindent}{0pt}
\usepackage{adjustbox}\DeclareUnicodeCharacter{2236}{:}\DeclareUnicodeCharacter{2190}{\ensuremath{\leftarrow}}
\providecommand{\R}{\mathbb{R}}
\providecommand{\N}{\mathbb{N}}
\providecommand{\Z}{\mathbb{Z}}
\providecommand{\Q}{\mathbb{Q}}
\providecommand{\D}{\mathbb{D}}
\providecommand{\C}{\mathbb{C}}
\providecommand{\e}{{\mathrm{e}}}
\providecommand{\vect}[1]{\overrightarrow{#1}}
\providecommand{\vectt}[1]{\overrightarrow{\mathrm{#1}}}
\providecommand{\vvt}[1]{\overrightarrow{\mathrm{#1}}}
\providecommand{\np}[2][]{{#2~#1}}
\providecommand{\vv}[1]{{\overrightarrow{#1}}}
\renewcommand{\d}{{\mathrm{d}}}
\providecommand{\degres}{{^\circ}}
\providecommand{\degre}{{^\circ}}
\providecommand{\blue}{\color{blue}}
\providecommand{\red}{\color{red}}
\providecommand{\enskip}{\quad}
\providecommand{\llbracket}{\lbrack\!\lbrack}
\providecommand{\rrbracket}{\rbrack\!\rbrack}
\providecommand{\euro}{{\text{€}}}
\providecommand{\no}{{\text{n°}}}
\providecommand{\interval}[2]{{[#1;#2]}}
\providecommand{\intervFO}[2]{{[#1;#2[}}
\providecommand{\dotfill}{\ldots}

\providecommand{\ds}{\displaystyle}
\providecommand{\textding}[1]{#1}
\providecommand{\Oij}{\(\mathrm{O};\vec{\imath},\vec{\jmath}\)}
\providecommand{\Oijk}{\(\mathrm{O};\vec{\imath},\vec{\jmath},\vec{k}\)}
\providecommand{\Ouv}{\(\mathrm{O};\vec{u},\vec{v}\)}
\providecommand{\intervFF}[2]{\ensuremath{[#1;#2]}}
\renewcommand{\euro}{\ifmmode\text{\texteuro}\else\texteuro\fi}
\DeclareUnicodeCharacter{00B0}{\ifmmode^{\circ}\else\textdegree\fi}
\DeclareUnicodeCharacter{20AC}{\ifmmode\text{\texteuro}\else\texteuro\fi}
% Énoncés : Éducation nationale. Transcription et corrigés : APMEP, Denis Vergès et contributeurs. Découpage et classement : CoopMaths, Sébastien Lozano, Éric Elter et contributeurs. Adaptation de la présentation : Maths Explorer.
% CC BY-SA : https://coopmaths.fr/www/about/
% https://coopmaths.fr/alea/static/bac/2023/tex/bac_2023_03_sujet2_ameriquenord_4.tex
\begin{document}
\noindent\textit{APMEP / CoopMaths. CC BY-SA. Présentation adaptée par Maths Explorer.}\par\medskip
\medskip

\emph{Cet exercice est un questionnaire à choix multiple. \\Pour chaque question, une seule des quatre réponses proposées est exacte.\\Le candidat indiquera sur sa copie le numéro de la question et la réponse choisie.\\Aucune justification n'est demandée.\\
Une réponse fausse, une réponse multiple ou l'absence de réponse à une question ne rapporte ni n'enlève de point.\\Les cinq questions sont indépendantes.}

\medskip

\begin{enumerate}
\item On considère la fonction $f$ définie sur l'intervalle $]1~;~+\infty[$ par 
\[f(x)= 0,05 - \dfrac{\ln x}{x - 1}.\]

La limite de la fonction $f$ en $+\infty$ est égale à :

\begin{center}
\begin{tabularx}{\linewidth}{*{4}{X}}
\textbf{a.~~} $+\infty$&\textbf{b.~~}0,05&\textbf{c.~~}$-\infty$&\textbf{d.~~}0
\end{tabularx}
\end{center}

\item On considère une fonction $h$ continue sur l'intervalle $[-2 ; 4]$ telle que :

\[h(-1)=0, \qquad h(1) = 4, \qquad h(3) = -1.\]

On peut affirmer que :

\begin{center}
\begin{tabularx}{\linewidth}{X}
\textbf{a.~~} la fonction $h$ est croissante sur l'intervalle $[-1~;~1]$.\\
\textbf{b.~~} la fonction $h$ est positive sur l'intervalle $[-1~;~1]$.\\
\textbf{c.~~} il existe au moins un nombre réel $a$ dans l'intervalle $[1~;~3]$ tel que $h(a) = 1$.\\
\textbf{d.~~} l'équation $h(x)=1$ admet exactement deux solutions dans l'intervalle 

$[-2~;~4]$.
\end{tabularx}
\end{center}

\item On considère deux suites $\left(u_{n}\right)$ et $\left(v_{n}\right)$ à termes strictement positifs telles que 

$\displaystyle\lim _{n \rightarrow+\infty} u_{n}=+\infty$ et $\left(v_{n}\right)$ converge vers 0 .

On peut affirmer que :

\begin{center}
\begin{tabularx}{\linewidth}{*{2}{X}}
\textbf{a.~~}la suite $\left(\frac{1}{v_{n}}\right)$ converge.&
\textbf{b.~~} la suite $\left(\frac{v_{n}}{u_{n}}\right)$ converge.\\
\textbf{c.~~} la suite $\left(u_{n}\right)$ est croissante.&
\textbf{d.~~} $\displaystyle\lim _{n \rightarrow+\infty}\left(-u_{n}\right)^{n}=-\infty$
\end{tabularx}
\end{center}

\item Pour participer à un jeu, un joueur doit payer $4$ \euro.

Il lance ensuite un dé équilibré à six faces :

\begin{itemize}
\item s'il obtient 1, il remporte $12$ \euro{};
\item s'il obtient un nombre pair, il remporte $3$\euro{};
\item sinon, il ne remporte rien.
\end{itemize}

En moyenne, le joueur :

\begin{center}
\begin{tabularx}{\linewidth}{*{2}{X}}
\textbf{a.~~}gagne $3,50$ \euro&\textbf{b.~~}perd $3$ \euro.\\
\textbf{c.~~}perd $1,50$ \euro&\textbf{d.~~}perd $0,50$ \euro.
\end{tabularx}
\end{center}

\item On considère la variable aléatoire $X$ suivant la loi binomiale $\mathcal{B}(3~;~p)$.

On sait que $P(X = 0) = \dfrac{1}{125}$. On peut affirmer que :

\begin{center}
\renewcommand\arraystretch{1.9}
\begin{tabularx}{\linewidth}{*{2}{X}}
\textbf{a.~~}$p = \dfrac{1}{5}$&\textbf{b.~~}$P(X = 1) =\dfrac{124}{125}$\\
\textbf{c.~~}$p = \dfrac{4}{5}$&\textbf{d.~~}$P(X= 1) =\dfrac{4}{5}$
\end{tabularx}
\end{center}
\end{enumerate}

\end{document}
